% \ThisChapterColor{A8DADC}   % optional: override the color from Colors.tex
\chapter{Introduction}
\label{ch:introduction}
\glsresetall

% (optional) an epigraph or picture for the colored opening page goes here
\clearpage   % end of the colored opening page

\section{Background}
\label{sec:intro:background}
\lipsum[1]

Acronyms are written once in full: \gls{udt}, and afterwards short: \gls{udt}.

\section{Research objectives}
\label{sec:intro:objectives}
The main research question of this dissertation is:

\mainrq

\begin{ResearchQuestion}[Short name]\label{rq:1}
    What is ... ?
\end{ResearchQuestion}

\begin{ResearchQuestion}[Short name]\label{rq:2}
    How can ... ?
\end{ResearchQuestion}

\section{Thesis outline}
\label{sec:intro:outline}
\Cref{ch:paper1} addresses \cref{rq:1}; \cref{ch:paper2} addresses \cref{rq:2};
\cref{ch:synthesis} synthesises the findings.

\begin{figure}[h]
    \centering
    % \includegraphics[width=\linewidth]{figures/thesis_outline}  % file in this chapter's figures/
    \fbox{\parbox[c][3cm][c]{0.8\linewidth}{\centering Thesis outline figure}}
    \caption{Structure of the dissertation.}
    \label{fig:intro:outline}
\end{figure}
